\section{Conclusion}
\label{sec:conclusion}

The $2n$ that a Frobenius-invariant factor base saves is a fact about
counts; collecting it is a question of representation. We have given
the representation, proved the three things an implementer needs to
know about it (orbits are full for prime $n$, the sign can be chosen
Frobenius-invariantly, and the least rotation of a necklace begins at
a longest zero run), measured what it costs where it is paid, and
shown with verified logarithms which of its consequences survive a
complete pipeline. The honest summary is that the representation is
cheap enough to be the right default on both sides of every
comparison, that on the index-calculus side it buys memory and
constants, and that on the Pollard rho side it buys a factor four in
the reference, which is why a fair comparison needs it on both sides.
The next measurements are fixed by the gaps in
Section~\ref{sec:limitations}: an isolated replay of the zero-gap stage
at $n=83$, an isolated one-target comparison at $n=83$, and a
branch-free two-limb canonicalizer for degrees no vector register
holds.
